% ============================================================
%  Chapter 5 — Arc Length and Asymptotic Equivalence
% ============================================================
\section{Arc Length and Asymptotic Equivalence}\label{sec:arclength}

In this section we derive the arc-length formula for the parabola $k(n)$
and prove that the arc length converges asymptotically to the
difference values $\Delta k$. The universal quantity $Q$ provides a
unified framework for comparing the arc lengths of the parabola and the
helix.

\subsection{Arc-Length Formula for the Parabola}

The arc-length element of the curve $(n, k(n))$ is
$ds = \sqrt{1 + [k'(n)]^2}\,dn$ with $k'(n) = (4n+3)/6$.
Substituting $u = 4n+3$, $du = 4\,dn$:

\begin{theorem}[Closed-form arc length]\label{thm:arc-length-parabola}
Let $u_i = 4n_i + 3$. The arc length of $k(n)$ between $n_1$ and $n_2$ is:
\begin{equation}\label{eq:arc-length-closed}
  L(n_1, n_2)
  = \frac{1}{48}\Bigl[u\sqrt{u^2+36}
    + 36\operatorname{arcsinh}\!\left(\frac{u}{6}\right)
    \Bigr]_{u_1}^{u_2}.
\end{equation}
\end{theorem}

\begin{proof}
From $ds = \frac{1}{6}\sqrt{36 + (4n+3)^2}\,dn$ and $u = 4n+3$:
$L = \frac{1}{24}\int_{u_1}^{u_2}\sqrt{u^2+36}\,du$.
Applying the standard integral
$\int\sqrt{u^2+a^2}\,du = \frac{1}{2}[u\sqrt{u^2+a^2} +
a^2\operatorname{arcsinh}(u/a)]$ with $a = 6$ and combining the
prefactors $\frac{1}{24}\cdot\frac{1}{2} = \frac{1}{48}$ yields the
result.
\end{proof}

\subsection{Unified \texorpdfstring{$Q$}{Q}-Representation}
\label{subsec:Q-arclength}

Since $u = 4n+3 = Q(k(n))$, the arc-length element can be written
directly in terms of $Q$:

\begin{theorem}[$Q$-parametrisation of the arc lengths]\label{thm:Q-arc}
The arc-length elements of the parabola and of the conical helix
admit the unified form
\begin{equation}\label{eq:ds-Q}
  ds = \frac{1}{24}\sqrt{Q^2 + C}\; dQ,
\end{equation}
where
\[
  C_{\mathrm{Parabola}} = 36,
  \qquad
  C_{\mathrm{Helix}} = 36 + \frac{144}{\pi^2} \approx 50.59.
\]
\end{theorem}

\begin{proof}
For the parabola, this is immediate from the substitution $u = Q$.
For the helix, using $\theta = \pi Q/12$ and $d\theta = (\pi/12)\,dQ$
in the speed formula~\eqref{eq:speed}:
\begin{align*}
  ds &= \|\gamma'(\theta)\|\,d\theta
     = \frac{3}{\pi^2}\sqrt{4(1+\theta^2)+\pi^2}\cdot\frac{\pi}{12}\,dQ \\
     &= \frac{1}{4\pi}\sqrt{4+4\theta^2+\pi^2}\,dQ
     = \frac{1}{24}\sqrt{Q^2 + 36 + \tfrac{144}{\pi^2}}\,dQ.
     \qedhere
\end{align*}
\end{proof}

\begin{remark}[Convergence of the arc lengths]
The difference $C_{\mathrm{Helix}} - C_{\mathrm{Parabola}} = 144/\pi^2$
is the exact contribution of the spiral motion. For $Q \to \infty$ the
ratio
$ds_{\mathrm{Helix}}/ds_{\mathrm{Parabola}} = \sqrt{(Q^2+C_H)/(Q^2+C_P)}
\to 1$, so the arc lengths of both curves converge.
\end{remark}

\subsection{Asymptotic Equivalence: Arc Length \texorpdfstring{$\approx \Delta k$}{≈ Δk}}

The most striking consequence of the arc-length formula is that the
difference values $\Delta k$ of the sequence are themselves excellent
approximations of the arc length.

\begin{proposition}[Asymptotic arc-length--difference equivalence]
\label{prop:asymptotic}
Let $L^{(2)}(m)$ and $L^{(4)}(m)$ denote the arc lengths of the
parabola over the $m$-th short step and the $m$-th long step,
respectively (Definition~\ref{def:steps}). Then, as $m \to \infty$:
\[
  \frac{L^{(2)}(m)}{\delta^{(2)}_m} \to 1,
  \qquad
  \frac{L^{(4)}(m)}{\delta^{(4)}_m} \to 1.
\]
\end{proposition}

\begin{proof}
We demonstrate the short-step case; the long-step case is analogous.
The short step connects $n_1 = 6m-1$ to $n_2 = 6m+1$, giving
$u_1 = 24m-1$ and $u_2 = 24m+7$.

By Theorem~\ref{thm:arc-length-parabola}, the arc length is
$L^{(2)}(m) = \frac{1}{48}[F(u)]_{u_1}^{u_2}$ where
$F(u) = u\sqrt{u^2+36} + 36\operatorname{arcsinh}(u/6)$.
For large $u$ the expansion
$\sqrt{u^2+36} = u + 18/u + O(u^{-3})$ gives:
\[
  F(u) = u^2 + 18 + 36\ln u + O(u^{-1}).
\]
Hence
\[
  F(u_2) - F(u_1)
  = (u_2^2 - u_1^2) + 36\ln\frac{u_2}{u_1} + O(m^{-1}).
\]
Now $u_2^2 - u_1^2 = (u_2+u_1)(u_2-u_1) = (48m+6)\cdot 8 = 384m + 48$,
so
\[
  L^{(2)}(m) = \frac{384m+48}{48} + O\!\left(\frac{\ln m}{m}\right)
             = 8m + 1 + O\!\left(\frac{\ln m}{m}\right)
             = \delta^{(2)}_m + o(1).
             \qedhere
\]
\end{proof}

\begin{table}[H]
\centering
\caption{Arc length vs.\ difference value. The ratio $L/\Delta k$
  converges monotonically to~$1$.}
\renewcommand{\arraystretch}{1.15}
\begin{tabular}{cccccc}
\toprule
$n_1 \to n_2$ & $\Delta n$ & $\Delta k$ & Arc length $L$ & $L / \Delta k$ \\
\midrule
 $1 \to  5$ & 4 & 10 & 10.839 & 1.0839 \\
 $5 \to  7$ & 2 &  9 &  9.221 & 1.0246 \\
 $7 \to 11$ & 4 & 26 & 26.310 & 1.0119 \\
$11 \to 13$ & 2 & 17 & 17.117 & 1.0069 \\
$13 \to 17$ & 4 & 42 & 42.191 & 1.0045 \\
$17 \to 19$ & 2 & 25 & 25.080 & 1.0032 \\
$19 \to 23$ & 4 & 58 & 58.138 & 1.0024 \\
$23 \to 25$ & 2 & 33 & 33.061 & 1.0018 \\
\bottomrule
\end{tabular}
\end{table}

\begin{remark}[Geometric meaning]
This convergence has a transparent cause: the slope
$k'(n) = (4n+3)/6 = Q/6$ grows without bound, so the horizontal
component $\Delta n \leq 4$ contributes ever less to the arc-length
element $\sqrt{1 + (k')^2}$. In the limit, $ds/dn \to k'(n)$ and
hence $L \to \int k'\,dn = \Delta k$.
\end{remark}

\subsection{Taylor Approximation}

The asymptotic equivalence $L \approx \Delta k$ can also be confirmed
through the $p$-parametrisation of the \emph{spiral} arc length
(cf.\ Section~\ref{sec:helix}), where the integrand takes the form
$\sqrt{1 + 3/(\pi^2 p)}$. Expanding for large $p$:

\begin{theorem}[Taylor approximation of the spiral arc length]
\label{thm:taylor}
For the spiral arc length in the $p$-parametrisation, with
$p_i = k(n_i)$:
\begin{equation}\label{eq:taylor}
  L \approx (p_2 - p_1)
  + \frac{3}{2\pi^2}\ln\!\left(\frac{p_2}{p_1}\right)
  + \frac{9}{8\pi^4}\left(\frac{1}{p_1} - \frac{1}{p_2}\right).
\end{equation}
The leading term $(p_2 - p_1) = \Delta k$ confirms the asymptotic
equivalence; the logarithmic correction is of relative order
$O(\ln p / p)$.
\end{theorem}
